% ============================================================================
% Soluciones del Módulo 5 - Bases de datos, SQL y ETL desde R
% ============================================================================

\begin{solucion}{m5-1}
Creamos la base en memoria, cargamos el CSV con \codigo{dbWriteTable()} y
exploramos. Primero los paquetes (en todas las soluciones de este módulo
usamos los mismos):

\begin{rcode}
library(tidyverse)
library(DBI)
library(RSQLite)
library(dbplyr)
\end{rcode}

\begin{rcode}
con <- dbConnect(RSQLite::SQLite(), ":memory:")

productos <- read_csv("datasets/productos.csv", show_col_types = FALSE)
dbWriteTable(con, "productos", productos)

dbListTables(con)
dbListFields(con, "productos")

dbGetQuery(con, "
  SELECT categoria,
         COUNT(*)                      AS productos,
         ROUND(AVG(precio_catalogo), 2) AS precio_promedio
  FROM productos
  GROUP BY categoria
  ORDER BY productos DESC, categoria
")

dbDisconnect(con)
\end{rcode}

\begin{rsalida}
[1] "productos"
 [1] "producto_id"     "nombre_producto" "categoria"      
 [4] "subcategoria"    "marca"           "precio_catalogo"
 [7] "costo"           "stock_actual"    "stock_minimo"   
[10] "proveedor"       "activo"          "margen_pct"     
[13] "estado_stock"   
     categoria productos precio_promedio
1      Belleza         9          474.54
2  Electrónica         7          413.72
3    Alimentos         6          354.92
4     Juguetes         6          425.89
5   Automotriz         5          365.35
6     Deportes         4          452.79
7       Libros         4          500.17
8     Mascotas         4          388.89
9        Hogar         3          361.51
10        Ropa         2          357.82
\end{rsalida}

La tabla tiene 13 columnas. Belleza es la categoría con más productos (9) y
Ropa la que menos (2). Libros tiene el precio de catálogo promedio más
alto (\$500.17), aunque con solo 4 productos.
\end{solucion}

\begin{solucion}{m5-2}
El filtro combina dos condiciones con \codigo{AND}; el faltante es una
columna calculada con alias, que podemos usar en el \codigo{ORDER BY}.

\begin{rcode}
con <- dbConnect(RSQLite::SQLite(), ":memory:")
dbWriteTable(con, "productos",
             read_csv("datasets/productos.csv", show_col_types = FALSE))

alerta <- dbGetQuery(con, "
  SELECT nombre_producto, categoria, proveedor, stock_actual,
         stock_minimo, stock_minimo - stock_actual AS faltante
  FROM productos
  WHERE activo = 1                       -- TRUE se guardó como 1
    AND stock_actual < stock_minimo
  ORDER BY faltante DESC
")
alerta

# ¿Qué proveedor aparece más?
dbGetQuery(con, "
  SELECT proveedor, COUNT(*) AS productos_en_alerta
  FROM productos
  WHERE activo = 1 AND stock_actual < stock_minimo
  GROUP BY proveedor
  ORDER BY productos_en_alerta DESC
")
dbDisconnect(con)
\end{rcode}

\begin{rsalida}
  nombre_producto  categoria   proveedor stock_actual stock_minimo
1     Producto 03 Automotriz Proveedor 3           22           26
2     Producto 33  Alimentos Proveedor 1           42           45
3     Producto 06   Juguetes Proveedor 1           36           38
  faltante
1        4
2        3
3        2
    proveedor productos_en_alerta
1 Proveedor 1                   2
2 Proveedor 3                   1
\end{rsalida}

Solo tres productos activos están por debajo de su mínimo, todos por pocas
piezas (el Producto 03 es el más urgente, con 4 faltantes). El Proveedor 1
surte dos de ellos: es con quien conviene hablar primero.
\end{solucion}

\begin{solucion}{m5-3}
\codigo{HAVING} filtra los grupos después de agregarlos.

\begin{rcode}
con <- dbConnect(RSQLite::SQLite(), ":memory:")
dbWriteTable(con, "empleados",
             read_csv("datasets/empleados.csv", show_col_types = FALSE) %>%
               mutate(fecha_ingreso = format(fecha_ingreso, "%Y-%m-%d")))

dbGetQuery(con, "
  SELECT departamento,
         COUNT(*)                       AS empleados,
         ROUND(AVG(salario_mensual), 2) AS salario_promedio,
         MAX(salario_mensual)           AS salario_maximo
  FROM empleados
  GROUP BY departamento
  HAVING COUNT(*) >= 15
  ORDER BY salario_promedio DESC
")
dbDisconnect(con)
\end{rcode}

\begin{rsalida}
         departamento empleados salario_promedio salario_maximo
1              Ventas        16         29031.76       43777.92
2 Atención al Cliente        29         27918.75       44438.77
3                RRHH        15         22955.96       42797.09
\end{rsalida}

Solo tres departamentos tienen 15 empleados o más. Ventas paga el salario
promedio más alto (\$29\,031.76) y Atención al Cliente es el departamento
más grande (29 personas).
\end{solucion}

\begin{solucion}{m5-4}
Reutilizamos la idea de \codigo{leer\_para\_bd()} del capítulo para
guardar las fechas como texto ISO. En la parte (b), la condición de
diciembre va \emph{dentro} del \codigo{ON} del \codigo{LEFT JOIN}: si la
pusiéramos en el \codigo{WHERE}, eliminaríamos precisamente los productos
sin pareja que buscamos.

\begin{rcode}
con <- dbConnect(RSQLite::SQLite(), ":memory:")
leer_para_bd <- function(archivo) {
  read_csv(file.path("datasets", archivo), show_col_types = FALSE) %>%
    mutate(across(where(is.Date), ~ format(.x, "%Y-%m-%d")),
           across(where(hms::is_hms), as.character))
}
for (tabla in c("transacciones", "tiendas", "productos")) {
  dbWriteTable(con, tabla, leer_para_bd(paste0(tabla, ".csv")))
}

# (a) Ventas por ciudad de la tienda
dbGetQuery(con, "
  SELECT ti.ciudad,
         COUNT(*)                           AS tickets,
         ROUND(SUM(t.total_transaccion), 0) AS ventas,
         ROUND(AVG(t.total_transaccion), 2) AS ticket_promedio
  FROM transacciones AS t
  INNER JOIN tiendas AS ti ON t.tienda_id = ti.tienda_id
  GROUP BY ti.ciudad
  ORDER BY ventas DESC
")

# (b) Productos sin ventas en diciembre
dbGetQuery(con, "
  SELECT p.producto_id, p.nombre_producto, p.categoria
  FROM productos AS p
  LEFT JOIN transacciones AS t
         ON p.producto_id = t.producto_id
        AND t.fecha BETWEEN '2023-12-01' AND '2023-12-31'
  WHERE t.transaccion_id IS NULL
  ORDER BY p.producto_id
")
dbDisconnect(con)
\end{rcode}

\begin{rsalida}
       ciudad tickets ventas ticket_promedio
1        CDMX     398 383588          963.79
2 Guadalajara     184 169260          919.89
3     Tijuana     106 114377         1079.03
4   Monterrey     109 105305          966.10
5   Querétaro     108 102791          951.77
6      Puebla      95  84672          891.29
   producto_id nombre_producto   categoria
1            5     Producto 05   Alimentos
2           10     Producto 10       Hogar
3           13     Producto 13        Ropa
4           18     Producto 18      Libros
5           23     Producto 23     Belleza
6           25     Producto 25    Juguetes
7           38     Producto 38   Alimentos
8           40     Producto 40 Electrónica
9           43     Producto 43  Automotriz
10          45     Producto 45 Electrónica
11          49     Producto 49     Belleza
\end{rsalida}

CDMX concentra las ventas (\$383\,588) porque ahí están cuatro de las diez
tiendas, pero Tijuana, con una sola tienda (el Outlet), tiene el ticket
promedio más alto (\$1\,079.03). En diciembre, 11 de los 50 productos no
tuvieron ninguna venta; si además tienen mucho inventario, son candidatos a
promoción.
\end{solucion}

\begin{solucion}{m5-5}
La tabla \codigo{ventas\_mal} guarda el número de días desde 1970 y
\codigo{strftime()} lo malinterpreta; la tabla \codigo{ventas} guarda
texto ISO y todo funciona.

\begin{rcode}
con <- dbConnect(RSQLite::SQLite(), ":memory:")
ventas_r <- read_csv("datasets/ventas_retail.csv", show_col_types = FALSE)

dbWriteTable(con, "ventas_mal", ventas_r)
dbWriteTable(con, "ventas",
             ventas_r %>% mutate(fecha = format(fecha, "%Y-%m-%d")))

dbGetQuery(con, "SELECT fecha, typeof(fecha) AS tipo,
                        strftime('%m', fecha) AS mes
                 FROM ventas_mal LIMIT 2")
dbGetQuery(con, "SELECT fecha, typeof(fecha) AS tipo,
                        strftime('%m', fecha) AS mes
                 FROM ventas LIMIT 2")

dbGetQuery(con, "
  SELECT (CAST(strftime('%m', fecha) AS INTEGER) + 2) / 3 AS trimestre,
         ROUND(SUM(total), 2)         AS ventas,
         ROUND(AVG(descuento_pct), 2) AS descuento_prom,
         SUM(descuento_pct > 0)       AS dias_con_descuento
  FROM ventas
  GROUP BY trimestre
  ORDER BY trimestre
")
dbDisconnect(con)
\end{rcode}

\begin{rsalida}
  fecha tipo mes
1 19358 real  11
2 19359 real  11
       fecha tipo mes
1 2023-01-01 text  01
2 2023-01-02 text  01
  trimestre   ventas descuento_prom dias_con_descuento
1         1 126273.9           4.33                 45
2         2 120716.9           4.51                 38
3         3 134427.4           5.43                 49
4         4 132498.3           4.84                 45
\end{rsalida}

En la tabla mal cargada, la fecha es un número (19358 son los días desde
1970 del 1 de enero de 2023) y \codigo{strftime()} devuelve el mes
equivocado (\codigo{11}) sin marcar error. En la tabla correcta, la fecha es
texto y el mes es \codigo{01}. El tercer trimestre fue el de más ventas
(\$134\,427.4) y también el de mayor descuento promedio (5.43\%).
\end{solucion}

\begin{solucion}{m5-6}
Los marcadores \codigo{:cliente}, \codigo{:desde} y \codigo{:hasta} se
llenan con la lista \codigo{params}. El valor nunca se interpreta como
SQL.

\begin{rcode}
con <- dbConnect(RSQLite::SQLite(), ":memory:")
dbWriteTable(con, "transacciones",
             read_csv("datasets/transacciones.csv",
                      show_col_types = FALSE) %>%
               mutate(fecha = format(fecha, "%Y-%m-%d"),
                      hora = as.character(hora)))

ventas_cliente <- function(con, cliente, desde, hasta) {
  dbGetQuery(con, "
    SELECT COUNT(*)                         AS compras,
           ROUND(SUM(total_transaccion), 2) AS gasto
    FROM transacciones
    WHERE cliente_id = :cliente
      AND fecha BETWEEN :desde AND :hasta",
    params = list(cliente = cliente,
                  desde   = format(as.Date(desde)),
                  hasta   = format(as.Date(hasta))))
}

ventas_cliente(con, 5, "2023-01-01", "2023-06-30")
ventas_cliente(con, 5, "2023-07-01", "2023-12-31")

# Intento de inyección: se busca un cliente_id igual al texto
ventas_cliente(con, "5 OR 1=1", "2023-01-01", "2023-12-31")

# Lo que habría pasado con paste0() (¡no lo hagas!)
dbGetQuery(con, paste0("SELECT COUNT(*) AS compras FROM transacciones ",
                       "WHERE cliente_id = ", "5 OR 1=1"))
dbDisconnect(con)
\end{rcode}

\begin{rsalida}
  compras   gasto
1       4 4560.68
  compras    gasto
1       9 10119.42
  compras gasto
1       0    NA
  compras
1    1000
\end{rsalida}

El cliente 5 hizo 4 compras por \$4\,560.68 en el primer semestre y 9 por
\$10\,119.42 en el segundo. Con el texto \codigo{"5 OR 1=1"} la consulta
parametrizada no encuentra nada (0 compras; la suma de cero filas es
\codigo{NA}), mientras que la versión con \codigo{paste0()} habría
devuelto las 1\,000 transacciones de todos los clientes.
\end{solucion}

\begin{solucion}{m5-7}
Con \paquete{dbplyr} construimos la consulta con verbos de dplyr; el
trimestre ya viene como columna en la tabla.

\begin{rcode}
library(dbplyr)
con <- dbConnect(RSQLite::SQLite(), ":memory:")
dbWriteTable(con, "transacciones",
             read_csv("datasets/transacciones.csv",
                      show_col_types = FALSE) %>%
               mutate(fecha = format(fecha, "%Y-%m-%d"),
                      hora = as.character(hora)))

consulta <- tbl(con, "transacciones") %>%
  group_by(metodo_pago, trimestre) %>%
  summarise(tickets = n(),
            ventas  = sum(total_transaccion, na.rm = TRUE),
            .groups = "drop") %>%
  arrange(metodo_pago, trimestre)

show_query(consulta)
resultado_dbplyr <- collect(consulta)
head(resultado_dbplyr, 4)

resultado_sql <- dbGetQuery(con, "
  SELECT metodo_pago, trimestre, COUNT(*) AS tickets,
         SUM(total_transaccion) AS ventas
  FROM transacciones
  GROUP BY metodo_pago, trimestre
  ORDER BY metodo_pago, trimestre")

all.equal(as.data.frame(resultado_dbplyr), resultado_sql)
dbDisconnect(con)
\end{rcode}

\begin{rsalida}
<SQL>
SELECT
  `metodo_pago`,
  `trimestre`,
  COUNT(*) AS `tickets`,
  SUM(`total_transaccion`) AS `ventas`
FROM `transacciones`
GROUP BY `metodo_pago`, `trimestre`
ORDER BY `metodo_pago`, `trimestre`
# A tibble: 4 x 4
  metodo_pago trimestre tickets ventas
  <chr>       <chr>       <int>  <dbl>
1 Efectivo    Q1             90 80565.
2 Efectivo    Q2             73 71206.
3 Efectivo    Q3             73 73187.
4 Efectivo    Q4             79 78777.
[1] TRUE
\end{rsalida}

El SQL generado es prácticamente idéntico al que escribirías a mano, y
\codigo{all.equal()} devuelve \codigo{TRUE}: ambos caminos dan el mismo
resultado. Convertimos el tibble con \codigo{as.data.frame()} porque
\codigo{dbGetQuery()} devuelve un data frame simple y \codigo{all.equal()}
también compara la clase.
\end{solucion}

\begin{solucion}{m5-8}
En (a) calculamos el gasto por cliente en una CTE y lo rankeamos dentro de
cada ciudad. En (b), \codigo{SUM(ventas) OVER (ORDER BY ventas DESC)} da el
acumulado y \codigo{SUM(ventas) OVER ()} (ventana vacía) el total de todas
las filas.

\begin{rcode}
con <- dbConnect(RSQLite::SQLite(), ":memory:")
dbWriteTable(con, "clientes",
             read_csv("datasets/clientes.csv", show_col_types = FALSE) %>%
               mutate(fecha_registro = format(fecha_registro)))
dbWriteTable(con, "transacciones",
             read_csv("datasets/transacciones.csv",
                      show_col_types = FALSE) %>%
               mutate(fecha = format(fecha), hora = as.character(hora)))

# (a) Los dos mejores clientes de cada ciudad
dbGetQuery(con, "
  WITH gasto AS (
    SELECT c.ciudad, c.nombre, SUM(t.total_transaccion) AS gasto
    FROM transacciones AS t
    JOIN clientes AS c ON t.cliente_id = c.cliente_id
    GROUP BY c.cliente_id
  ),
  ranking AS (
    SELECT ciudad, nombre, ROUND(gasto, 0) AS gasto,
           RANK() OVER (PARTITION BY ciudad ORDER BY gasto DESC) AS lugar
    FROM gasto
  )
  SELECT * FROM ranking
  WHERE lugar <= 2 AND ciudad IN ('Monterrey', 'Guadalajara', 'CDMX')
  ORDER BY ciudad, lugar
")

# (b) Curva de Pareto de productos
pareto <- dbGetQuery(con, "
  WITH ventas_producto AS (
    SELECT producto_id, SUM(total_transaccion) AS ventas
    FROM transacciones
    GROUP BY producto_id
  )
  SELECT producto_id, ROUND(ventas, 0) AS ventas,
         ROW_NUMBER() OVER (ORDER BY ventas DESC) AS posicion,
         ROUND(100.0 * SUM(ventas) OVER (ORDER BY ventas DESC) /
               SUM(ventas) OVER (), 1) AS pct_acumulado
  FROM ventas_producto
  ORDER BY ventas DESC
")
head(pareto, 3)
pareto %>% filter(pct_acumulado >= 80) %>% head(1)
dbDisconnect(con)
\end{rcode}

\begin{rsalida}
       ciudad             nombre gasto lugar
1        CDMX        Carmen Cruz 10438     1
2        CDMX Fernando Hernández  8242     2
3 Guadalajara     Carmen Jiménez 11848     1
4 Guadalajara      Alberto Reyes 11356     2
5   Monterrey      Marta Ramírez  6261     1
6   Monterrey   Isabel Hernández  5888     2
  producto_id ventas posicion pct_acumulado
1           6  34028        1           3.5
2          19  28368        2           6.5
3          31  25430        3           9.1
  producto_id ventas posicion pct_acumulado
1          22  16949       37          81.4
\end{rsalida}

En Guadalajara los dos mejores clientes gastaron más de \$11\,000 cada uno,
mientras que en Monterrey el mejor apenas llega a \$6\,261. En la curva de
Pareto, el 80\% de las ventas se alcanza hasta el producto en la posición
37: se necesitan 37 de 50 productos (74\%). Las ventas están muy
repartidas, lejos del clásico ``20\% de los productos genera el 80\% de las
ventas''.
\end{solucion}

\begin{solucion}{m5-9}
Primero el DDL con llaves y reglas; luego la carga con
\codigo{dbAppendTable()} (que respeta la estructura), la vista y el
índice. La fecha es la llave primaria de los hechos porque hay una venta
por día.

\begin{rcode}
con <- dbConnect(RSQLite::SQLite(), ":memory:")
dbExecute(con, "PRAGMA foreign_keys = ON")
dbExecute(con, "CREATE TABLE dim_producto (
  producto_id INTEGER PRIMARY KEY, nombre_producto TEXT, categoria TEXT)")
dbExecute(con, "CREATE TABLE dim_tienda (
  tienda_id INTEGER PRIMARY KEY, nombre_tienda TEXT, tipo TEXT)")
dbExecute(con, "CREATE TABLE fact_ventas_diarias (
  fecha TEXT PRIMARY KEY,
  producto_id INTEGER NOT NULL REFERENCES dim_producto(producto_id),
  tienda_id   INTEGER NOT NULL REFERENCES dim_tienda(tienda_id),
  cantidad    INTEGER CHECK (cantidad > 0),
  total REAL)")

dbAppendTable(con, "dim_producto",
  read_csv("datasets/productos.csv", show_col_types = FALSE) %>%
    transmute(producto_id = as.integer(producto_id), nombre_producto,
              categoria))
dbAppendTable(con, "dim_tienda",
  read_csv("datasets/tiendas.csv", show_col_types = FALSE) %>%
    transmute(tienda_id = as.integer(tienda_id), nombre_tienda, tipo))
dbAppendTable(con, "fact_ventas_diarias",
  read_csv("datasets/ventas_retail.csv", show_col_types = FALSE) %>%
    transmute(fecha = format(fecha, "%Y-%m-%d"),
              producto_id = as.integer(producto_id),
              tienda_id = as.integer(tienda_id),
              cantidad = as.integer(cantidad), total))

dbExecute(con, "CREATE VIEW vw_trimestre_tipo AS
  SELECT (CAST(strftime('%m', f.fecha) AS INTEGER) + 2) / 3 AS trimestre,
         t.tipo, COUNT(*) AS dias, ROUND(SUM(f.total), 0) AS ventas
  FROM fact_ventas_diarias AS f
  JOIN dim_tienda AS t ON f.tienda_id = t.tienda_id
  GROUP BY trimestre, t.tipo")
dbGetQuery(con, "SELECT * FROM vw_trimestre_tipo
                 WHERE trimestre = 1 ORDER BY ventas DESC")

consulta <- "SELECT SUM(total) FROM fact_ventas_diarias WHERE tienda_id = 3"
dbGetQuery(con, paste("EXPLAIN QUERY PLAN", consulta))$detail
dbExecute(con, "CREATE INDEX idx_tienda ON fact_ventas_diarias(tienda_id)")
dbGetQuery(con, paste("EXPLAIN QUERY PLAN", consulta))$detail
dbDisconnect(con)
\end{rcode}

\begin{rsalida}
[1] 0
[1] 0
[1] 0
[1] 0
[1] 50
[1] 10
[1] 365
[1] 0
  trimestre      tipo dias ventas
1         1  Estándar   46  54838
2         1 Principal   14  30699
3         1    Outlet   14  17226
4         1   Premium   10  16706
5         1   Express    6   6804
[1] "SCAN fact_ventas_diarias"
[1] 0
[1] "SEARCH fact_ventas_diarias USING INDEX idx_tienda (tienda_id=?)"
\end{rsalida}

Los \codigo{[1] 0} son las instrucciones DDL y los \codigo{[1] 50},
\codigo{10} y \codigo{365} las filas cargadas por \codigo{dbAppendTable()}.
En el primer trimestre, las tiendas de tipo Estándar suman más ventas
porque son seis. El plan pasa de \codigo{SCAN} (recorrer toda la tabla) a
\codigo{SEARCH \ldots{} USING INDEX idx\_tienda} tras crear el índice.
\end{solucion}

\begin{solucion}{m5-10}
La transformación \emph{no} filtra cantidades negativas a propósito: la
regla \codigo{CHECK} de la base es la última línea de defensa y queremos
ver cómo actúa. La carga usa una tabla temporal y un \codigo{INSERT} que
solo toma las fechas que aún no existen. ¿Por qué no \codigo{INSERT OR
IGNORE}? Porque en SQLite \codigo{OR IGNORE} ignora en silencio
\emph{cualquier} violación de restricción, no solo las llaves repetidas:
también la de \codigo{CHECK}. El lote con cantidad negativa se
``cargaría'' con estado OK y 0 filas nuevas, sin ningún error. Si lo
pruebas, lo verás.

\begin{rcode}
con <- dbConnect(RSQLite::SQLite(), ":memory:")
dbExecute(con, "CREATE TABLE fact_ventas_diarias (
  fecha TEXT PRIMARY KEY, producto_id INTEGER, cliente_id INTEGER,
  tienda_id INTEGER, cantidad INTEGER CHECK (cantidad > 0),
  precio_unitario REAL, descuento_pct REAL, total REAL)")
dbExecute(con, "CREATE TABLE log_etl (
  id INTEGER PRIMARY KEY AUTOINCREMENT, fecha_hora TEXT, lote TEXT,
  leidas INTEGER, nuevas INTEGER, estado TEXT, mensaje TEXT)")

extraer_ventas <- function(ruta) read_csv(ruta, show_col_types = FALSE)

transformar_ventas <- function(datos) {
  datos %>%
    distinct(fecha, .keep_all = TRUE) %>%
    transmute(fecha = format(as.Date(fecha), "%Y-%m-%d"),
              producto_id = as.integer(producto_id),
              cliente_id = as.integer(cliente_id),
              tienda_id = as.integer(tienda_id),
              cantidad = as.integer(cantidad), precio_unitario,
              descuento_pct, total)
}

cargar_ventas <- function(con, datos) {
  dbWriteTable(con, "stg", datos, temporary = TRUE, overwrite = TRUE)
  # Solo las fechas que aún no existen (sin OR IGNORE: ver explicación)
  n <- dbExecute(con, "INSERT INTO fact_ventas_diarias
                       SELECT * FROM stg
                       WHERE fecha NOT IN (SELECT fecha
                                           FROM fact_ventas_diarias)")
  dbExecute(con, "DROP TABLE stg")
  n
}

ejecutar_etl <- function(con, ruta) {
  leidas <- NA_integer_
  en_transaccion <- FALSE
  registrar <- function(nuevas, estado, mensaje = NA_character_) {
    dbAppendTable(con, "log_etl", tibble(
      fecha_hora = format(Sys.time()), lote = basename(ruta),
      leidas = leidas, nuevas = nuevas, estado = estado,
      mensaje = mensaje))
  }
  tryCatch({
    datos <- extraer_ventas(ruta)
    leidas <- nrow(datos)
    limpios <- transformar_ventas(datos)
    dbBegin(con); en_transaccion <- TRUE
    nuevas <- cargar_ventas(con, limpios)
    registrar(nuevas, "OK")
    dbCommit(con)
  }, error = function(e) {
    if (en_transaccion) dbRollback(con)
    registrar(0L, "ERROR", conditionMessage(e))
  })
  invisible(NULL)
}

# Simular los lotes en una carpeta temporal
carpeta <- file.path(tempdir(), "lotes")
dir.create(carpeta, showWarnings = FALSE)
ventas <- read_csv("datasets/ventas_retail.csv", show_col_types = FALSE)
lote <- function(desde, hasta, nombre) {
  ruta <- file.path(carpeta, nombre)
  ventas %>% filter(fecha >= as.Date(desde), fecha <= as.Date(hasta)) %>%
    write_csv(ruta)
  ruta
}
l1 <- lote("2023-01-01", "2023-06-30", "lote_ene_jun.csv")
l2 <- lote("2023-04-01", "2023-09-30", "lote_abr_sep.csv")
l3 <- lote("2023-10-01", "2023-12-31", "lote_oct_dic.csv")
l_malo <- file.path(carpeta, "lote_malo.csv")
tibble(fecha = as.Date("2024-01-01"), producto_id = 1, cliente_id = 1,
       tienda_id = 1, cantidad = -2, precio_unitario = 100,
       descuento_pct = 0, total = -200) %>% write_csv(l_malo)

for (ruta in c(l1, l2, l3, l3, l_malo)) ejecutar_etl(con, ruta)

dbGetQuery(con, "SELECT id, lote, leidas, nuevas, estado FROM log_etl")
dbGetQuery(con, "SELECT mensaje FROM log_etl WHERE estado = 'ERROR'")

# Validación final contra el CSV
dbGetQuery(con, "SELECT COUNT(*) AS filas, ROUND(SUM(total), 2) AS total
                 FROM fact_ventas_diarias")
round(sum(ventas$total), 2)
dbDisconnect(con)
\end{rcode}

\begin{rsalida}
[1] 0
[1] 0
  id             lote leidas nuevas estado
1  1 lote_ene_jun.csv    181    181     OK
2  2 lote_abr_sep.csv    183     92     OK
3  3 lote_oct_dic.csv     92     92     OK
4  4 lote_oct_dic.csv     92      0     OK
5  5    lote_malo.csv      1      0  ERROR
                                mensaje
1 CHECK constraint failed: cantidad > 0
  filas    total
1   365 513916.6
[1] 513916.6
\end{rsalida}

La bitácora cuenta la historia: el segundo lote leyó 183 días pero solo
cargó 92 (abril-junio ya estaban), el lote repetido no cargó nada y el
lote malo terminó en \codigo{ERROR} porque la regla \codigo{CHECK} rechazó
la cantidad negativa. La tabla final tiene exactamente 365 filas y el mismo
total que el CSV: el ETL es idempotente, atómico y trazable.
\end{solucion}
